\section{Experiments}
\label{experiments}

The experiment section evaluates whether the proposed workflow improves environment-aware requirement elicitation.
No result values are reported in this draft structure.

\subsection{Research Questions}
\label{exp:rqs}

\begin{itemize}[leftmargin=*]
  \item \textbf{RQ1:} Does the workflow improve requirement context completeness?
  \item \textbf{RQ2:} Does knowledge-assisted analysis improve environment consistency?
  \item \textbf{RQ3:} Do dynamic clarification questions help stakeholders complete missing environmental information?
  \item \textbf{RQ4:} What is the contribution of each agent role and knowledge component?
\end{itemize}

\subsection{Experimental Setup}
\label{exp:setup}

This subsection will describe the elicitation tasks, participant or annotator protocol, model configuration, prompt setting, and evaluation procedure.
The setup should focus on environment-aware elicitation rather than general text generation quality.

\subsection{Datasets}
\label{exp:datasets}

The dataset should consist of natural language requirement descriptions with missing or underspecified environmental context.
Candidate sources may include course project requirements, curated requirement scenarios, and manually designed environment-sensitive cases.

\subsection{Baselines}
\label{exp:baselines}

Baselines should be simple and directly comparable:
\begin{itemize}[leftmargin=*]
  \item original requirements without environment-aware elicitation;
  \item single-agent requirement clarification;
  \item multi-agent workflow without environment knowledge;
  \item workflow without human clarification, if applicable.
\end{itemize}

\subsection{Evaluation Metrics}
\label{exp:metrics}

\subsubsection{Context Completeness}
\label{exp:context_completeness}

Measures whether deployment, device, runtime, and user-context information is sufficiently specified.

\subsubsection{Environment Consistency}
\label{exp:environment_consistency}

Measures whether elicited environmental assumptions are mutually consistent and aligned with the scenario.

\subsubsection{Requirement Completeness}
\label{exp:requirement_completeness}

Measures whether the resulting requirements contain enough information for environment-aware analysis.

\subsubsection{Clarification Effectiveness}
\label{exp:clarification_effectiveness}

Measures whether generated clarification questions help stakeholders provide missing environmental information.

\subsection{Main Results}
\label{exp:main_results}

This subsection will report quantitative and qualitative findings for the research questions after experiments are conducted.
Tables should distinguish confirmed results from illustrative examples.

\subsection{Ablation Study}
\label{exp:ablation}

The ablation study will compare the full workflow against variants without selected agent roles, knowledge assistance, or human clarification.

\subsection{Case Study}
\label{exp:case_study}

The case study will trace one requirement from its initial underspecified form through environment detection, clarification, and context-enriched requirement output.
